% !TEX program = pdflatex
% !TEX encoding = UTF-8 Unicode

% Plantilla de la clase `scrbook` del paquete KOMA-script para la
% elaboración de un TFG siguiendo las directrices del la comisión del
% Grado en Matemáticas de la Universidad de Granada.

% Francisco Torralbo Torralbo
% miércoles, 29 de abril de 2020

\documentclass{scrbook}

\KOMAoptions{%
  fontsize=10pt,        % Tamaño de fuente
  paper=a4,             % Tamaño del papel
  headings=normal,      % Tamaño de letra para los títulos: small, normal, big
  % parskip=half,         % Espacio entre párrafos: full (una línea) o half (media línea)
  headsepline=false,    % Una linea separa la cabecera del texto
  cleardoublepage=empty,% No imprime cabecera ni pie en páginas en blanco
  chapterprefix=false,  % No antepone el texto "capítulo" antes del número
  appendixprefix=false,	% No antepone el texto "Apéndice" antes de la letra
  listof=totoc,		    	% Añade a la tabla de contenidos la lista de tablas y figuras
  index=totoc,			    % Añade a la talba de contenidos una entrada para el índice
  bibliography=totoc,	  % Añade a la tabla de contenidos una entrada para bibliografía
  BCOR=5mm,           % Reserva margen interior para la encuadernación.
                        % El valor dependerá el tipo de encuadernado y del grosor del libro.
  DIV=10,             % Cálcula el diseño de página según ciertos
                        % parámetros. Al aumentar el número aumentamos el ancho de texto y disminuimos el ancho del margen. Una opción de 14 producirá márgenes estrechos y texto ancho.
}

% INFORMACIÓN PARA LA VERSIÓN IMPRESA
% Si el documento ha de ser impreso en papel de tamaño a4 pero el tamaño del documento (elegido en \KOMAoptions con la ocpión paper) no es a4 descomentar la línea que carga el paquete `crop` más abajo. El paquete crop se encargará de centrar el documento en un a4 e imprimir unas guías de corte. El procedimiento completo para imprenta sería el siguiente:
% 0. Determinar, según el tipo de encuadernación del documento, el ancho reservado para el proceso de encuadernación (preguntar en la imprenta), es decir, la anchura del área del papel que se pierde durante el proceso de encuadernación. Fijar la varibale BCOR de \KOMAoptions a dicho valor.
% 1. Descomentar la siguiente línea e imprimir una única página con las guías de corte
% 2. Cambiar la opción `cross` por `cam` (o `off`) en el paquete crop y volver a compilar. Imprimir el documento (las guías de corte impresas no inferfieren con el texto).
% 3. Usar la página con las guías impresas en el punto 1 para cortar todas las páginas.

% \usepackage[a4, odd, center, pdflatex, cross]{crop} % Permite imprimir el documento en un a4 (si el tamaño es más pequeño) mostrando unas guías de corte. Útil para imprenta.

% VERSIÓN ELECTRÓNICA PARA TABLETA
% Las opciones siguientes seleccionan un tamaño de impresión similar a una tableta de 9 pulgadas con márgenes estrechos. Útil para producir una versión en pdf para ser leída en una tableta en lugar de impresa.
% Para que la portada quede centrada correctamente hay que editar el
% archivo `portada.tex` y eliminar el entorno `addmargin`

% \KOMAoptions{fontsize=10pt, paper=19.7104cm:14.7828cm, twoside=false, BCOR=0cm, DIV=14}

% ---------------------------------------------------------------------
%	PAQUETES
% ---------------------------------------------------------------------

% CODIFICACIÓN E IDIOMA
% ---------------------------------------------------------------------
\usepackage[utf8]{inputenc} 			    % Codificación de caracteres

% Selección del idioma: cargamos por defecto inglés y español (aunque este último es el idioma por defecto para el documento). Cuando queramos cambiar de idioma escribiremos:
% \selectlanguage{english} o \selectlanguage{spanish}

\usepackage[english, spanish, es-nodecimaldot, es-noindentfirst, es-tabla]{babel}

% Opciones cargadas para el paquete babel:
  % es-nodecimaldot: No cambia el punto decimal por una coma en modo matemático.
  % es-noindentfirst: No sangra los párrafos tras los títulos.
  % es-tabla: cambia el título del entorno `table` de "Cuadro" a "Tabla"

% Otras opciones del paquete spanish-babel:
  \unaccentedoperators % Desactiva los acentos en los operadores matemáticso (p.e. \lim, \max, ...). Eliminar esta opción si queremos que vayan acentuados

% MATEMÁTICAS
% ---------------------------------------------------------------------
\usepackage{amsmath, amsthm, amssymb} % Paquetes matemáticas
\usepackage{mathtools}                % Añade mejoras a amsmath
\mathtoolsset{showonlyrefs=true}      % sólo se numeran las ecuaciones que se usan
\usepackage[mathscr]{eucal} 					% Proporciona el comando \mathscr para
                                      % fuentes de tipo manuscrito en modo matemático sin sobreescribir el comando \mathcal

% TIPOGRAFÍA
% ---------------------------------------------------------------------
% El paquete microtype mejora la tipografía del documento.
\usepackage[activate={true,nocompatibility},final,tracking=true,kerning=true,spacing=true,factor=1100,stretch=10,shrink=10]{microtype}

% Las tipografías elegidas para el documento son las siguientes
% Normal font: 			URW Palladio typeface.
% Sans-serif font: 	Iwona
% Monospace font: 	Inconsolata
% Consultar http://www.tug.dk/FontCatalogue/ para seleccionar otra tipografía.
% Es conveniente elegir aquellas que tienen soporte matemático.
\usepackage[T1]{fontenc}
\usepackage[sc, osf]{mathpazo} \linespread{1.05}
\usepackage[scaled=.95,type1]{cabin} % sans serif in style of Gill Sans
\usepackage{inconsolata}
% \renewcommand{\sfdefault}{iwona}

% Selecciona el tipo de fuente para los títulos (capítulo, sección, subsección) del documento.
\setkomafont{disposition}{\sffamily\bfseries}

% Cambia el ancho de la cita. Al inicio de un capítulo podemos usar el comando \dictum[autor]{cita} para añadir una cita famosa de un autor.
\renewcommand{\dictumwidth}{0.45\textwidth}

\recalctypearea % Necesario tras definir la tipografía a usar.

% TABLAS, GRÁFICOS, LISTADOS DE CÓDIGO y ALGORITMOS
% ---------------------------------------------------------------------
\usepackage{booktabs}
% \renewcommand{\arraystretch}{1.5} % Aumenta el espacio vertical entre las filas de un entorno tabular

\usepackage{xcolor, graphicx}
% Carpeta donde buscar los archivos de imagen por defecto
\graphicspath{{img/}}

% IMAGEN DE LA PORTADA
% Existen varias opciones para la imagen de fondo de la portada del TFG. Todas ellas tienen en logotipo de la universidad de Granada en la cabecera. Las opciones son las siguientes:
% 1. portada-ugr y portada-ugr-color: diseño con marca de agua basada en el logo de la UGR (en escala de grises y color).
% 2. portada-ugr-sencilla y portada-ugr-sencilla-color: portada únicamente con el logotipo de la UGR en la cabecera.
\usepackage{eso-pic}
\newcommand\BackgroundPic{%
	\put(0,0){%
		\parbox[b][\paperheight]{\paperwidth}{%
			\vfill
			\centering
      % Indicar la imagen de fondo en el siguiente comando
			\includegraphics[width=\paperwidth,height=\paperheight,%
			keepaspectratio]{portada-ugr-sencilla}%
			\vfill
}}}

\usepackage{listings} % Para la inclusión de trozos de código

% Configuración del código
\lstset{
	basicstyle=\ttfamily,%
	breaklines=true,%
	captionpos=b,                    % sets the caption-position to bottom
  tabsize=2,	                   % sets default tabsize to 2 spaces
  frame=none,
  numbers=left,
  xleftmargin=18pt,
  stepnumber=1,
  aboveskip=12pt,
  showstringspaces=false,
  keywordstyle=\bfseries,
  commentstyle=\itshape,
  numberstyle=\scriptsize\bfseries,
  morekeywords={sage},
}

\usepackage[spanish,ruled,vlined]{algorithm2e} % Para la inclusión de algoritmos

% CABECERAS
% ---------------------------------------------------------------------
% Si queremos modificar las cabeceras del documento podemos usar el paquete
% `scrlayer-scrpage` de KOMA-Script. Consultar la documentación al respecto.
% \usepackage[automark]{scrlayer-scrpage}

% VARIOS
% ---------------------------------------------------------------------

%\usepackage{showkeys}	% Muestra las etiquetas del documento. Útil para revisar las referencias cruzadas.

% ÍNDICE
% Para generar el índice hay que compilar el documento con MakeIndex. Generalmente los editores se encargan de ello automáticamente.
% ----------------------------------------------------------------------
% \index{} para añadir un elemento
% \index{main!sub} para añadir un elementos "sub" bajo la categoría "main".
% \index{termino|textbf} para dar formato al número de página (negrita).
% \index{termino|see{termino relacionado}} para crear una referencia cruzada

% Ejemplo: \index{espacio homogéneo}, \index{superficie!mínima}, \index{esfera|see{espacio homogéneo}}

\usepackage{makeidx}
%\usepackage{showidx} % Muestra en el margen del documento las entradas añadidas al índice. Útil para revisar el documento. Si está activo el índice no se genera
\makeindex

% ---------------------------------------------------------------------
% COMANDOS Y ENTORNOS
% ---------------------------------------------------------------------
% Cargamos un archivo externo donde hemos incluido todos los comandos
% propios que vamos a usar en el documento.
% DEFINICIÓN DE COMANDOS Y ENTORNOS

% CONJUNTOS DE NÚMEROS

  \newcommand{\N}{\mathbb{N}}     % Naturales
  \newcommand{\Z}{\mathbb{Z}}     % Enteros
  \newcommand{\Q}{\mathbb{Q}}     % Racionales
  \newcommand{\C}{\mathcal{C}}    % Código
  \newcommand{\R}{\mathcal{R}}    % Álgebra Cociente
  \newcommand{\F}{\mathbb{F}}     % Cuerpo

% Letras
  \newcommand{\w}{\omega}

% Operadores
  \newcommand{\lc}{\operatorname{lc}}       % Coeficiente Lider
  \newcommand{\prob}{\operatorname{prob}}   % Probabilidad

% TEOREMAS Y ENTORNOS ASOCIADOS

  \theoremstyle{plain}
  \newtheorem{theorem}{Teorema}
  \newtheorem{proposition}[theorem]{Proposición}
  \newtheorem{lemma}[theorem]{Lema}
  \newtheorem{corollary}[theorem]{Corolario}
  \theoremstyle{definition}
  \newtheorem{definition}[theorem]{Definición}
  \newtheorem{proofs}{Demostración}
  \theoremstyle{remark}
  \newtheorem{remark}[theorem]{Observación}
  \newtheorem{exampleth}[theorem]{Ejemplo}


% --------------------------------------------------------------------
% INFORMACIÓN DEL TFG Y EL AUTOR
% --------------------------------------------------------------------
\usepackage{xspace} % Para problemas de espaciado al definir comandos

\newcommand{\miTitulo}{Algoritmo de Sugiyama para códigos skew RS\xspace}
\newcommand{\miNombre}{Guillermo Galindo Ortuño\xspace}
\newcommand{\miGrado}{Doble Grado en Ingeniería Informática y Matemáticas}
\newcommand{\miFacultad}{Escuela Técnica Superior de Ingenierías Informática y de Telecomunicación
\\Facultad de Ciencias}
\newcommand{\miUniversidad}{ Universidad de Granada}
% Añadir tantos tutores como sea necesario separando cada uno de ellos
% mediante el comando `\\\medskip` y una línea en blanco
\newcommand{\miTutor}{
  José Gómez Torrecillas \\ \emph{Departamento de Álgebra}
  \\\medskip

  Francisco Javier Lobillo Borrego \\ \emph{Departamento de Álgebra}
}
\newcommand{\miCurso}{2019-2020\xspace}

% HYPERREFERENCES
% --------------------------------------------------------------------
\usepackage{xurl}
\usepackage[pagebackref]{hyperref}
% Opciones para el paquete hyperref
%----------------------------------

\hypersetup{%
  % hidelinks,            % Enlaces sin color ni borde. El borde no se imprime
  linkbordercolor=0.8 0 0,
  citebordercolor=0 0.8 0,
  citebordercolor=0 0.8 0,
  colorlinks = true,            % Color en texto de los enlaces. Comentar esta línea o cambiar `true` por `false` para imprimir el documento.
  linkcolor = [rgb]{0.5, 0, 0}, % Color de los enlaces internos
  urlcolor = [rgb]{0, 0, 0.5},  % Color de los hipervínculos
  citecolor = [rgb]{0, 0.5, 0}, % Color de las referencias bibliográficas
	pdftitle={\miTitulo},%
	pdfauthor={\textcopyright\ \miNombre, \miFacultad, \miUniversidad},%
  pdfsubject={Trabajo de fin de Grado},%
	pdfkeywords={},%
	pdfcreator={pdfLaTeX},%
}

% Redefinición del estilo para mostrar las referencias cruzadas en la bibliografía.
\renewcommand*{\backref}[1]{}
\renewcommand*{\backrefalt}[4]{{\footnotesize [%
    \ifcase #1 No citado%
    \or Citado en pág.~#2%
    \else Citado en págs. #2%
    \fi%
]}}

% Etiquetas en español para el comando \autoref
\def\chapterautorefname{Capítulo}
\def\appendixautorefname{Apéndice}
\def\sectionautorefname{Sección}
\def\subsectionautorefname{Subsección}
\def\figureautorefname{Fig.}
\def\tableautorefname{Tabla}

\def\teoremaautorefname{Teorema}
\def\proposicionautorefname{Proposición}
\def\lemaautorefname{Lema}
\def\corolarioautorefname{Corolario}
\def\definicionautorefname{Def.}
\def\observacionautorefname{Observación}
\def\ejemploautorefname{E.j.}

% Pone automáticamente un parántesis para las ecuaciones
\def\equationautorefname~#1\null{Ec.~(#1)\null}


\begin{document}

% --------------------------------------------------------------------
% FRONTMATTER
% -------------------------------------------------------------------
\frontmatter % Desactiva la numeración de capítulos y usa numeración romana para las páginas

% \pagestyle{plain} % No imprime cabeceras

% !TeX root = ../libro.tex
% !TeX encoding = utf8

%*******************************************************
% Titlepage
%*******************************************************
\begin{titlepage}
  \AddToShipoutPicture*{\BackgroundPic}
  \phantomsection 
  \pdfbookmark[1]{Título}{title}

  % Para que el título esté centrado en la página.
  % Los valores numéricos deberán elegirse de acuerdo con el diseño de
  % página (sobre todo si se cambia la opción BCOR o DIV).
  \begin{addmargin}[2.575cm]{0cm}
  \begin{flushleft}
    \Large  
    \hfill\vfil

    \large{\textsf{\miFacultad}}
    \vfill

    {\large\textsc\miGrado} \vfill


    {\large\textsc{trabajo de fin de grado}}

    \begingroup
    \Huge{\miTitulo}
    \endgroup

    \vfill\vfill\vfill\vfill

    \textsf{\normalsize{Presentado por:}}\\
    {\normalsize\textrm{\miNombre}} 
    \bigskip

    \textsf{\normalsize{Tutor:}}\\
    {\normalsize\rmfamily \miTutor}

    \bigskip
    \textsf{\normalsize{Curso académico \miCurso}}
  \end{flushleft}  
  \end{addmargin}       

\end{titlepage}   
\cleardoublepage
\endinput

% !TeX root = ../libro.tex
% !TeX encoding = utf8

%*******************************************************
% Little Dirty Titlepage
%*******************************************************

\thispagestyle{empty}

\begin{center}
  \large  

  \vspace*{\stretch{1}}

  \begingroup
  \huge{\miTitulo} \\ \bigskip
  \endgroup

  \textrm{\miNombre}

  \vspace{\stretch{5}}

\end{center}  

\newpage
\thispagestyle{empty}

\hfill

\vfill

\miNombre \textit{\miTitulo}.

Trabajo de fin de Grado. Curso académico \miCurso.
\bigskip

\begin{minipage}[t]{0.25\textwidth}
  \flushleft
  \textbf{Responsable de tutorización}
\end{minipage}
\begin{minipage}[t]{0.45\textwidth}
  \flushleft
  \miTutor
\end{minipage}
\begin{minipage}[t]{0.30\textwidth}
  \flushright
  \miGrado
  \medskip

  \miFacultad
  \medskip

  \miUniversidad
\end{minipage}

\newpage
\endinput

% !TeX root = ../libro.tex
% !TeX encoding = utf8
%
%*******************************************************
% Declaración de originalidad
%*******************************************************

\thispagestyle{empty}

\hfill\vfill

\textsc{Declaración de originalidad}\\\bigskip

D./Dña. \miNombre \\\medskip

Declaro explícitamente que el trabajo presentado como Trabajo de Fin de Grado (TFG), correspondiente al curso académico \miCurso, es original, entendida esta, en el sentido de que no ha utilizado para la elaboración del trabajo fuentes sin citarlas debidamente.
\medskip

En Granada a \today 
\begin{flushleft} 
Fdo: \miNombre 

\end{flushleft}

\vfill

\cleardoublepage
\endinput

% % !TeX root = ../libro.tex
% !TeX encoding = utf8

%*******************************************************
% Dedication
%*******************************************************
\thispagestyle{empty}
\phantomsection 
\pdfbookmark[1]{Dedicatoria}{Dedicatoria}

\hfill
\vfill

\begin{flushright}
\itshape
Dedicatoria (opcional) \\
Ver archivo \texttt{preliminares/dedicatoria.tex}
\end{flushright}

\vfill

\cleardoublepage
\endinput
                % Opcional
% !TeX root = ../libro.tex
% !TeX encoding = utf8

%*******************************************************
% Table of Contents
%*******************************************************
\phantomsection
\pdfbookmark[0]{\contentsname}{toc}

\setcounter{tocdepth}{2} % <-- 2 includes up to subsections in the ToC
\setcounter{secnumdepth}{3} % <-- 3 numbers up to subsubsections

% \manualmark
% \markboth{\textsc{\contentsname}}{\textsc{\contentsname}}
\tableofcontents

%*******************************************************
% List of Figures and of the Tables
%*******************************************************

%    % *******************************************************
%    %  List of Figures
%    % *******************************************************
%    \phantomsection
%    \listoffigures

%    %*******************************************************
%    % List of Tables
%    %*******************************************************
%    \phantomsection
%    \listoftables

%    %*******************************************************
%    % List of Listings
%    % The package \usepackage{listings} is needed
%    %*******************************************************
%	  % \phantomsection
%    % \renewcommand{\lstlistlistingname}{Listados de código}
%    % \lstlistoflistings

\cleardoublepage

% % !TeX root = ../libro.tex
% !TeX encoding = utf8

%*******************************************************
% Agradecimientos
%*******************************************************

\chapter{Agradecimientos}

En primer lugar, agredecer a mis dos tutores, José Gómez Torrecillas y Francisco
Javier Lobillo Borrego por su gran ayuda, consejos y completa disposición a lo
largo del desarrollo de todo este trabajo.

Quiero agredecer también a mi familia, a mis padres, Guillermo Galindo Maqueda
Ruiz y María Isabel Ortuño Ruiz, y a mi hermanos, Luis Galindo Ortuño e Isabel
Galindo Ortuño por todo el apoyo, la confianza, y el cariño recibidos durante todos
estos años, y sin los cuales no habría podido llegar hasta aquí.

A mis amigos, tanto de la residencia como compañeros de clase, que han hecho
que mi paso por Granada sea algo que no voy a olvidar, y que han estado presentes
tanto en los mejores como en los no tan buenos momentos de estos cinco años.

Por último, agredecer a toda la comunidad que contribuye en SageMath, por hacer posible
la implementación de este trabajo, y por su aportación al conocimiento científico y al
sofware de forma libre y gratuita.
\cleardoublepage
\endinput
            % Opcional

% \pagestyle{scrheadings} % A partir de ahora sí imprime cabeceras

% !TeX root = ../libro.tex
% !TeX encoding = utf8
%
%*******************************************************
% Summary
%*******************************************************

\selectlanguage{english}
\chapter{Summary}

The main goal of this project is to present, study and implement a version of the Sugiyama algoritm for a class of skew cyclic codes, the skew RS codes. In order to accomplish this, firstly we need to study the basics of linear codes and Ore extensions, which we cover in the first and second chapter of this thesis.\\

In the first chapter, we introduce the basic concepts about linear codes. In essence, a linear code over a finite field is a vector subspace with an associated distance, so it is a well-known mathematical structure. One of the most important concepts that is presented in this chapter is the generator matrix, which is the main tool for enconding messages into the code. The most common distance and the one that we study is the Hamming distance, defined as the number of coordinates that differ between two vectors. Next in this chapter, we discuss the decoding procedure, presenting two of the simpliest decoding algorithms and a proof of a result that stablishes the maximum number of errors that a linear code can correct. This latter result is crucial in coding theory, as research in this area tries to find efficient algorithms which can decode as many errors as possible. Finally, we briefly introduce cyclic codes, in which the codes we will construct in following chapters are based.\\

In the second chapter, we introduce Ore extensions. These have a similar structure as common polynomial rings over a field, with the difference that the product operation is not commutative. Instead, this product is \textit{twisted} by a field automorphism. The usual concepts about polynomial rings, as degree or leading coefficient, also exist for Ore extensions, and are included in this chapter. We will also introduce algorithms to calculate the division from left or right. These algorithms are essential and allow us to show that Ore extensions are non-commutative principal ideal domains, and to develop a left and right version of the extended Euclidean algorithm. The fact that the product is not commutative does not permit to think of the evaluation of these polynomials in the usual manner. Thus, there exist evaluations on the right and on the left, which are defined in accordance to the polynomial remainder theorem, and the respective division algorithm. In order to compute this evaluation, we present the notion of \textit{jth-norm}, and show its relation with the remainder of left and right division by linear factors. \\

In the third chapter, we will present the algorithm whose analysis and implementation in the SageMath framework is the main goal of this thesis, and with it all the mathematical structures and results needed. Along this whole chapter, \(\F[x;\sigma]\) will represent the Ore extension over \(\F\) given by the  \(\F\)-automorphism \(\sigma\). Then, we will firstly introduce the class of \textit{skew cyclic codes}, which are described as left ideals of the quotient ring \(\F[x;\sigma] / \langle x^{n} - 1 \rangle\) where \(n\) is the order of \(\sigma\). Therefore, the code itself is given by the image via the coordinate map (with respect to the usual basis) of the left ideal mentioned. Consecutively, we will present the necessary results to construct skew cyclic codes of a designed Hamming distance, which we denote by \textit{skew RS codes}. Later in this chapter, the concepts of \textit{error locator polynomial} and \textit{error evaluator polynomial} will be described, and we will show how they allow us to find and correct the errors in a received message by simply solving a linear system. Then we will prove a relation between these two polynomials, called the \textit{key equation}. Using this relation, we will be able to present the decoding algorithm and prove that it correctly decodes the message. Lastly, we will cover the procedure to follow when a \textit{key equation error} occurs in order to find the error locator and evaluator polynomials, and analyze how probable is the ocurrence of this type of errors.

In the fourth and last chapter, we present the documentation of the classes that we have developed for SageMath. This classes store the information needed to represent both skew cyclic codes and skew RS codes, and allow us to work with the already existing structures in SageMath. The classes implemented are:

\begin{itemize}
    \item A class to represent skew cyclic codes.
    \item Two classes to encode messages into skew cyclic code words, one for the vector form and another for the polynomial form.
    \item A class to represent skew RS codes.
    \item A class with the implementation of the decoding algorithm that is the goal of this thesis.
\end{itemize}

% Al finalizar el resumen en inglés, volvemos a seleccionar el idioma español para el documento
\selectlanguage{spanish}
\endinput

% !TeX root = ../libro.tex
% !TeX encoding = utf8
%
%*******************************************************
% Introducción
%*******************************************************

% \manualmark
% \markboth{\textsc{Introducción}}{\textsc{Introducción}} 

\chapter{Introducción}

De acuerdo con la comisión de grado, el TFG debe incluir una introducción en la que se describan claramente los objetivos previstos inicialmente en la propuesta de TFG, indicando si han sido o no alcanzados, los antecedentes importantes para el desarrollo, los resultados obtenidos, en su caso y las principales fuentes consultadas.

Ver archivo \texttt{preliminares/introduccion.tex}

\endinput


% --------------------------------------------------------------------
% MAINMATTER
% --------------------------------------------------------------------
\mainmatter % activa la numeración de capítulos, resetea la numeración de las páginas y usa números arábigos

\setpartpreamble[c][0.75\linewidth]{%
	\bigskip % Deja un espacio vertical en la parte superior
  Si el trabajo se divide en diferentes partes es posible incluir al inicio de cada una de ellas un breve resumen que indique el contenido de la misma. Esto es opcional.
}
% \part{Primera parte} % Dividir un libro en partes OPCIONAL

\chapter{Conceptos básicos sobre códigos lineales}%
\label{chap:conceptos_básicos_sobre_códigos_lineales}

\section{Códigos lineales, matrices generatriz y de paridad}%
\label{sec:códigos_lineales_matrices_generatriz_y_de_paridad}

Sea \(\F_q^n\) el espacio vectorial de todas las \(n\)-tuplas sobre el cuerpo finito \(\F_q\).

\begin{definition}[Código]
 Definimos un \((n, M)\) \textit{código} \(\mathcal{C}\) sobre \(\F_q\) como un subconjunto de \(\F_q^n\) de tamaño \(M\).
\end{definition}

Normalmente escribiremos los vectores \((a_1, a_2, \dots, a_n)\) en \(\F_q^n\) de la forma \(a_1 a_2 \cdots a_n\) y llamaremos a los vectores en \(\mathcal{C}\) \textit{palabras código}, o simplemente \textit{palabras}. Además, utilizaremos nombres concretos para referirnos a códigos sobre algunos de los cuerpos más comunes . A los códigos sobre \(\F_2\) los llamaremos \textit{códigos binarios}, los códigos sobre \(\F_3\) los notaremos como \textit{códigos ternarios} y a los códigos sobre \(F_4\) los llamaremos \textit{códigos cuaternarios}.\\

Dicho esto, necesitamos dotar de algún tipo de estructura adicional a los códigos pues siendo únicamente conjuntos estamos muy limitados. La estructura adicional más útil que se les impone es la linealidad.

\begin{definition}[Código lineal]
Si \(\mathcal{C}\) es un subespacio vectorial de dimensión \(k\) de \(\F_q^n\), diremos que \(\mathcal{C}\) es un \([n,k]\) \textit{código lineal} sobre \(\F_q\).
\end{definition}

Dicho esto, las dos maneras más comunes de presentar un código lineal son dando una matriz generatriz o una matriz de paridad.

\begin{definition}[Matriz generatriz]
    Una \textit{matriz generatriz} de un \([n,k]\) código lineal \(\mathcal{C}\) es cualquier matriz \(G\) de dimensiones \(k \times n\) cuyas filas formen una base de \(\mathcal{C}\).
\end{definition}

Dado una matriz generatriz \(G\), para cualquier conjunto de \(k\) columnas independientes de esta, diremos que el correspondiente conjunto de coordenadas es un \textit{conjunto de información} de C. Las restantes \(r=n-k\) coordenadas las notaremos como \textit{conjunto de redundancia}, y llamaremos a \(r\) \textit{redundancia} de \(\C\). Si las primeras \(k\) coordenadas forman un conjunto de información, existe una única matriz generatriz para el código de la forma \([I_k | A]\) donde \(I_k\) es la matriz identidad de orden \(k\). Diremos que una matriz generatriz así está en \textit{forma estándar}.

\begin{definition}
    Una \textit{matriz de paridad } de un \([n,k]\) código lineal \(\C\) es cualquier matriz \(H\) de dimensiones \((n-k)\times n\) tal que
    \[
    \C = \{x \in \F_q^n | Hx^T = 0\}
    .\]
\end{definition}

Como un código lineal es un subespacio de un espacio vectorial, es el núcleo de alguna aplicación lineal, y por tanto, para un código lineal siempre existe alguna matriz de paridad \(H\). Mencionemos que las filas de \(H\) son también independientes. Esto es porque, al ser \(H\) una aplicación lineal de \(\F_q^n\) en  \(\F_q^{n-k}\), y la dimensión del núcleo de dicha aplicación es \(k\), tenemos que la dimensión de la imagen es \(n-k\) y por tanto el rango de \(H\) también.

\section{Códigos duales}%
\label{sec:códigos_duales}

Como mencionamos en la sección anterior, las filas de una matriz de paridad de un código \(\C\) son independientes, y por tanto \(H\) es a su vez la matriz generatriz de otro código, el llamado \textit{dual} o \textit{ortogonal} de \(\C\), denotado por \(\C^\bot\). Es inmediato ver que \(\C^\bot\) es un \([n, n-k]\) código. Además, es posible dar una construcción alternativa para \(\C^\bot\) utilizando el producto escalar en  \(\F_q^n\):
\[
\C^\bot = \{x \in \F_q^n | x \cdot c = 0\ \forall c \in \C\}
.\]

\section{Pesos y distancias}
Una de las características más útiles para el estudio de códigos es la distancia mínima entre las palabras de este.

\begin{definition}
Dados dos vectores \(x, y \in F_q^n\), definimos la distancia \textit{Hamming} entre ellos \(d(x,y)\) como el numero de coordenadas en las que  \(x\) e \(y\) difieren.
\end{definition}

Veamos que en efecto esta es una distancia:

\begin{proposition}
La función distancia \(d(x,y)\) satisface las siguientes condiciones:
\begin{nlist}
    \item \(d(x,y) \geq 0\) para todo \(x, y \in \F_q^n\).
    \item \(d(x,y) = 0\) si y solo si \(x = y\).
    \item \(d(x,y) = d(y,x)\) for all \(x,y \in \F_q^n\)
    \item \(d(x,z) \leq d(x,y) + d(y,z)\) para todo \(x, y, z \in \F_q^n\)
\end{nlist}

\begin{proofs}
Las tres primeras propiedades son obvias por la propia definición de la distancia. Veamos pues la propiedad iv).\\

Dados dos vectores \(x, y \in \F_q^n\), definimos el conjunto \(D(x,y) = \{i | x_i \neq y_i\}\), y denotamos el complementario por \(D^c(x, y) = \{i | x_i = y_i\}\). Es claro que entonces el cardinal de \(D(x,y)\) coincide con nuestra distancia. \\

Recordemos también algunas propiedades sobre cardinales de conjuntos. Sea un conjunto \(A\), notemos por \(|A|\) su cardinal. Entonces, para cualesquiera conjuntos \(A\) y \(B\):
\begin{nlist}
    \item \(|A| \leq |A \cup B|\)
    \item \(|A \cup |B| \leq |A| + |B|\)
    \item Si \(|A| \leq |B|\), entonces \(|A^c| \geq |B^c|\)
\end{nlist}
Con esto, dados \(x, y, z \in \F_q^n\), conjuntos tenemos que

\[
D^c(x,z) = \{i | x_i = z_i\} = \{i | x_i = z_i = y_i\} \cup \{i | x_i = z_i \neq y_i\}
\]
\[
\implies |D^c(x,z)| \geq |\{i | x_i = z_i = y_i\}| = |\{i | x_i = y_i\} \cap \{i | z_i = y_i\}|
\]
\[
\implies |D(x,z)| \leq |\{i | x_i \neq y_i\} \cup \{i | z_i \neq y_i\}| = |D(x,y) \cup D(y,z)|
\]
\[
\implies |D(x,z)| \leq |D(x,y)| + |D(y,z)|
\]
quedándo demostrado el resultado.
\end{proofs}
\end{proposition}

Ahora, la \textit{distancia (mínima)} de un código \(\C\) es la mínima distancia entre dos palabras distintas de dicho código. Esta propiedad será crucial a la hora de determinar el número de errores que podrá corregir un código.

\begin{definition}
Diremos que el peso (\textit{de Hamming}) \(wt(x)\) de un vector \(x \in \F_q^n\) es el número de coordenadas distintas de cero de \(x\).
\end{definition}

Si tenemos dos vectores \(x, y \in \F_q^n\) es inmediato comprobar que \(d(x,y) = wt(x - y)\). En el siguiente mostramos la relación entre la distancia y el peso.

\begin{proposition}
Si \(\C\) es un  código lineal, la distancia mínima coincide con el mínimo de los pesos de las palabras distintas de cero de C.

\begin{proofs}
Sea \(d = d(x,y)\) la distancia mínima del código \(\C\), que se alcanza entre dos vectores \(x, y \in \C\), y \(d' = wt(z)\) el peso mínimo que se alcanza en \(z \in C\).

Por ser \(\C\) un subespacio vectorial,  \(0 \in \C\), y por tanto  \(d \leq d(z, 0) = wt(z) = d'\). De nuevo por ser \(\C\) un subespacio vectorial tenemos que \(x-y \in \C\), y por tanto  \(d' \leq wt(x-y) = d(x,y) = d\), y por tanto \(d = d'\).
\end{proofs}
\end{proposition}

Como consecuencia a este resultado, para códigos lineales, a la distancia mínima también se le llama \textit{peso mínimo} del código. En adelante, si el peso mínimo \(d\) de un \([n,k]\) código es conocido, entonces nos referiremos al código como un \([n,k,d]\) código.

\section{Codificar, decodificar, y el Teorema de Shannon}

\subsection{Codificar}%
Sea \(\C\) un \([n,k]\) código lineal sobre el cuerpo \(\F_q\) con matriz generatriz \(G\). Como este es un subespacio vectorial de \(\F_q^n\) de dimensión \(k\), contiene \(q^k\) palabras, que están en correspondencia uno a uno con \(q^k\) posibles mensajes. Por esto, la forma más simple es ver estos mensajes como \(k\)-tuplas \(x\) en \(\F_q^k\). Así, lo más común es codificar un mensaje \(x\) como la palabra \(c = xG\). Si G está en forma estándar, las primeras \(k\) coordenadas son los símbolos de información \(x\); el resto de \(n-k\) símbolos son los simbolos de paridad, es decir, la redundancia añadida a \(x\) con el fin de poder recuperarla si ocurre algún error. Dicho esto, la matriz \(G\) puede no estar en forma estándar. En particular, si existen índices de columnas \(i_1, i_2, \dots, i_n \) tales que la matriz \(k \times k\) formada por estas columnas es la matriz identidad, entonces el mensaje se encuentra en las coordenadas \(i_1, i_2, \dots, i_n \) separado pero sin modificar, es decir, el símbolo del mensaje \(x_j\) se encuentra en la componente \(i_j\) de la palabra código. Si esto ocurre diremos que el codificador e s \textit{sistemático}.

\subsection{Decodificar y el Teorema de Shannon}%
\label{sub:decodificar_y_el_teorema_de_shannon}

El proceso de decodificar, consistente en determinar qué palabra (y por tanto qué mensaje \(x\)) fue mandado al recibir un vector \(y\), es más complejo. Encontrar algoritmos de decodificación eficientes es una área de investigación muy relevante en la teoría de códigos debido a sus aplicaciones prácticas. En general, codificar es sencillo y decodificar es complicado, especialmente si tiene un tamaño suficientemente grande.

\chapter{Introducción a extensiones de Ore}%
\label{chap:conceptos_básicos_sobre_extensiones_ore}

A continuación introduciremos los conceptos principales sobre extensiones de Ore que utilizaremos como base para el siguiente capítulo. Las definiciones y principales resultados de este capítulo siguen el desarrollo realizado en~\cite{factoring_ore}.

\section{Conceptos básicos sobre extensiones de Ore}

Aunque las definiciones que realizamos a continuación se pueden hacer más generales, en este caso nos restringiremos al caso más particular necesario para continuar. Dicho esto, nuestro algoritmo trabajará sobre polinomios de Ore no conmutativos, con una única indeterminada \(x\), con coeficientes en un cuerpo cualquiera. Precisando, nuestros polinomios serán elementos de un anillo asociativo \(R = \F[x;\sigma]\), donde

\begin{itemize}
    \item \(\F\) es un cuerpo cualquiera.
    \item \(\sigma: \F \to \F\) es un automorfismo de cuerpos de orden finito digamos \(n\).
\end{itemize}

La construcción de \(R = \F[x;\sigma]\) sigue de la siguiente forma:
\begin{itemize}
    \item \(R\) es un \(\F\)-espacio vectorial a la izquierda sobre la base  \(\{x^n: n \geq 0\}\). Entonces, los elementos de  \(R\) son polinomios a la izquierda de la forma \(a_0 + a_1 x + \cdots + a_n x^n\) con \(a_i \in \F\).
    \item La suma de polimios es la usual.
    \item El producto de \(R\) está basado en las siguientes reglas: \(x^n x^m = x^{n+m}\), para \(m, n \in \mathbb{N}\), y \(xa = \sigma(a)x\) para  \(a \in \F\). Este producto se extiende recursivamente a \(R\).
\end{itemize}

Ahora que hemos definido formalmente las extensiones de Ore, introduciremos varios conceptos análogos a los anillos de polinomios conmutativos comunes.

El grado \(\deg (f)\) de un polinomio no nulo \(f \in R\), al igual que su coeficiente líder, se definen de la manera usual, de forma que
\[
f = \lc(f)x^{\deg(f)} + f_{\downarrow}
,\]
donde \(f_{\downarrow}\) es un polinomio de grado menor que \(\deg(f)\), y el coeficiente líder  \(\operatorname{lc}(f)\) es no nulo. Siguiendo las convenciones usuales para el polinomio nulo diremos que \(\deg(0) = -\infty\), y que \(\lc(0) = 0\).

Al igual que ocurre para polinomios conmutativos, directamente de la definición del producto de polinomios, dados \(f, g \in R\):
\[
\deg(fg) = \deg(f) + \deg(g)
.\]

Además, de la misma definición obtenemos también que
\[
\lc(fg) = \lc(f)\sigma^{\deg(f)}(\lc(g))
,\]

y por tanto \(R\) es un dominio de integridad no conmutativo.

El anillo \(R\) tiene algoritmos de división a la izquierda y derecha (veánse los algoritmos~\ref{alg:left_euclidean_div} y~\ref{alg:right_euclidean_div}).

\begin{algorithm}[H]
 \label{alg:left_euclidean_div}
 \SetKwInput{KwIn}{Entrada}
 \SetKwInput{KwOut}{Salida}
 \SetKw{Initialization}{Inicialización:}
 \KwIn{\(f,g \in \F[x;\sigma] \text{ con } g \neq 0\)}
 \KwOut{\(q,r \in \F[x;\sigma] \text{ tales que } f = qg + r \text{ y } \deg(r) < \deg(g)\)}
 \Initialization{\(q:=0, r:=f\)} \\
 \While{\(\deg(g) \leq \deg(r)\)}{
  \(a = \lc(r)\sigma^{deg(r) - deg(g)}(\lc(g)^{-1})\) \\
  \(q := q + ax^{\deg(r) - \deg(g)}, r:= r - ax^{deg(r) - deg(g)}g\)
 }
 \caption{División Euclídea a la izquierda}
\end{algorithm}

\begin{algorithm}[H]
 \label{alg:right_euclidean_div}
 \SetKwInput{KwIn}{Entrada}
 \SetKwInput{KwOut}{Salida}
 \SetKw{Initialization}{Inicialización:}
 \KwIn{\(f,g \in \F[x;\sigma] \text{ con } g \neq 0\)}
 \KwOut{\(q,r \in \F[x;\sigma] \text{ tales que } f = gq + r \text{ y } \deg(r) < \deg(g)\)}
 \Initialization{\(q:=0, r:=f\)} \\
 \While{\(\deg(g) \leq \deg(r)\)}{
  \(a = \sigma^{-deg(g)}(\lc(g)^{-1}\lc(r))\) \\
  \(q := q + ax^{\deg(r) - \deg(g)}, r:= r - gax^{deg(r) - deg(g)}\)
 }
 \caption{División Euclídea a la derecha}
\end{algorithm}

Mostramos a continuación una demostración que justifica estos algoritmos, cuya versión original puede encontrarse en~\cite[Th. 4.34]{bueso2003algorithmic}

\begin{theorem}
Sea \(\F\) un cuerpo finito de \(q\) elementos siendo \(q\) una potencia de un primo,  \(\sigma\) un autormorfismo de \(\F\) no nulo, y \(R = \F[x;\sigma]\) la extensión de Ore correspondiente. Entonces, dados \(f, g \in R\) existen \(q, r \in R\) únicos tales que:
\begin{enumerate}
    \item \(f = qg + r\).
    \item \(\deg(r) < \deg(g)\).
\end{enumerate}

Bajo las mismas hipótesis, existen también \(q, r \in R\) únicos tales que:
\begin{enumerate}
    \item \(f = gq + r\)
    \item \(\deg(r) < \deg(g)\).
\end{enumerate}
\end{theorem}

\begin{proof}
Para abreviar digamos \(m = \deg(g), n = \deg(f)\). Veamos primero la prueba de la división a la izquierda. Si \(m > n\), entonces no tenemos nada que probar, pues tomando  \(q = 0, r = f\) se cumple el resultado. Por otro lado, si \(m \leq n\), sean \(f = \sum_{i=0}^{n} a_i x^i\) y \(g = \sum_{j=0}^{m} b_j x^j\), aplicaremos inducción sobre \(n\). Si \(n = 0\), entonces también \(m = 0\), así que  \(f = a_0, g = b_0\), y por tanto tomamos \(r = 0, q = a_0 b_0^{-1}\).

Por tanto, supongamos la afirmación cierta para todo \(f\) de grado menor que \(n\).  Sea \(a = a_n \sigma^{n-m}(b_m^{-1})\). Entonces es claro que
\[
\deg(ax^{n-m}g) = n,
\]
\[
\lc(ax^{n-m}g) = a_n
.\]
Por tanto tenemos que
\[
\deg(f - a x^{n-m}g) < n,
\]
y por tanto, la hipótesis de inducción nos dice que existen \(q'\) y \(r'\) compliendo que \(\deg(r') < \deg(g)\)  y
\[
f - a x^{n-m}g = q'g + r'
.\]

Sea
\[
q = a x^{n-m} + q'
,\]

entonces
\[
f = a x^{n-m}g + q'g + r' = qg + r'
.\]

Queda probar que \(q\) y \(r\) son únicos como tales. Supongamos que
\[
f = q_1g + r_1 = q_2g + r_2,
\]
con \(\deg(r_1), \deg(r_2) < \deg(g)\). Entonces, \((q_1 - q_2)g = r_2 - r_1\), y podemos afirmar entonces que
\[
\deg(q_1 - q_2) + \deg(g) = \deg((q_1-q_2)g)
\]
\[
= \deg(r_2-r_1) \leq \max(\deg(r_2), \deg(r_1)) < \deg(g)
.\]
Esto prueba que \(\deg(q_1-q_2) = -\infty\), demostrando que \(q_1 - q_2 = 0\) y que \(r_2 = r_1 = 0\) que termina la prueba de la primera parte.

La prueba de la división a la derecha es completamente análoga, tomando \(a = \sigma^{-m}(a_{n}b_{m}^{-1})\), y utilizando que \(\deg(f - gax^{n-m}) < n\).

\end{proof}

Los polinomios \(r\) y \(q\) obtenidos como salida del algoritmo \ref{alg:left_euclidean_div} los llamaremos \textit{resto a la izquierda} y  \textit{cociente a la izquierda}, respectivamente, de la división a la izquierda de \(f\) por \(g\). Utilizaremos la notación \(r = \operatorname{lrem}(f,g)\) y \(q = \operatorname{lquo}(f,g)\). Asumimos convenciones y notaciones análogas para el algoritmo de división a la derecha.

Un polinomio \(f \in R\) se dice \textit{central} si para cualquier otro polinomio \(g \in R\), se tiene que \(fg = gf\).

\begin{lemma}
\label{lem:central_decomposition}
    Sea \(f \in R\) un polinomio central y \(g,h \in R\) tales que \(f = gh\). Entonces también se cumple la igualdad \(f = hg\)
\end{lemma}

\begin{proof}
    Multiplicando \(f\) a la derecha por \(g\) y usando que \(f\) es central tenemos que, \( fg = ghg = ggh\), y por tanto \(hg = gh = f.\)
\end{proof}

\section{Máximo común divisor y mínimo común multiplo}%
\label{sec:máximo_común_divisor_y_mínimo_común_multiplo}

Veamos ahora que, como consecuencia del algoritmo de división a la izquierda, dado un ideal a la izquierda \(I\) de \(R\), y cualquier polinomio no nulo  \(f \in I\) de grado mínimo, \(f\) es un generador de  \(I\). Notaremos en este caso \(I = Rf\).

En efecto, para cualquier \(g \in I\), utilizando el algoritmo~\ref{alg:left_euclidean_div} tenemos que \(g = q*f + r\) con \(\deg(r) < \deg(f)\). Pero \(g\) y \(q*f\) pertenecen a \(I\), y por tanto  \(r\) también. Como \(f\) era de grado mínimo en \(I\),  \(r = 0\).

Análogamente se prueba que cualquier ideal a la derecha de \(R\) es principal.  Por tanto \(R\) es un dominio de ideales principales no conmutativo.

Dados \(f,g \in R\), \(Rf \subset Rg\) significa que \(g\) es un \textit{divisor a la derecha} de \(f\), simbólicamente \(g|_{r} f\), o que \(f\) es \textit{múltiplo a la izquierda} de \(g\).

Por ser \(R\) un dominio de ideales principales sabemos que \(Rf + Rg = Rd\) para algún \(d \in R\), y es inmediato comprobar que \(d |_r f\) y  \(d |_r g\). Además, si tenemos \(d'\) con \(d' |_r f\), \(d' |_r g\), entonces \(Rf + Rg \subset Rd'\), luego \(Rd \subset Rd'\) y por tanto  \(d |_r d'\). En este caso diremos que \(d\) es un \textit{máximo común divisor a la derecha} de \(f\) y \(g\), estando unívocamente determinado salvo múltiplicación a la izquierda por una unidad de \(R\). Utilizaremos la notación \(d = {(f,g)}_r\). Además de aquí obtenemos directamente la \textbf{identidad de Bezout}.

\begin{proposition}[Identidad de bezout]
Sean \(f \text{ y } g\) dos polinomios en \(R\), y  \(d = {( f,g )_r}\), entonces existen \(u, v \in R\) tales que
\[
uf + vg = d
.\]
\end{proposition}

Similarmente  \(Rf \cap Rg = Rm\) si y solo si \(m\) es un \textit{mínimo común múltiplo a la izquierda} de \(f\) y \(g\), notado por \(m = {[f,g]}_l\). Este también es único salvo multiplicación a la izquierda por una unidad de \(R\).

La versión a la derecha de todas estas definiciones y propiedades puede establecerse de forma completamente análoga.

A continuación mostramos las respectivas versiones del Algoritmo Extendido de Euclides a derecha e izquierda (algoritmos~\ref{alg:right_ext_euc_alg} y~\ref{alg:left_ext_euc_alg} respectivamente). Estas nos permiten calcular el máximo común divisor y el mínimo común multiplo tanto a izquierda como a derecha.
En particular, para nuestro algoritmo principal utilizaremos la versión a la derecha de este algoritmo para obtener los coeficientes de Bezout.

\begin{algorithm}[H]\label{alg:right_ext_euc_alg}
 \SetKwInput{KwIn}{Entrada}
 \SetKwInput{KwOut}{Salida}
 \SetKw{Initialization}{Inicialización:}
 \KwIn{\(f,g \in \F[x;\sigma] \text{ con } f \neq 0,\ g \neq 0\)}
 \KwOut{\(\{u_i, v_i, r_i\}_{i = 0, \dots, h, h+1} \text{ tales que } r_i = fu_i + gv_i \text{ para todo  }\) \(0 \le i \le h+1, r_h = (f,g)_l, \text{ y } fu_{h+1} = [f,g]_r.\)}
 \Initialization{\\
 \(r_0 \leftarrow f, r_1 \leftarrow g\) \\
 \(u_0 \leftarrow 1, u_1 \leftarrow 0\) \\
 \(v_0 \leftarrow 0, v_1 \leftarrow 1\) \\
 \(q \leftarrow 0, rem \leftarrow 0\) \\
 \(i \leftarrow 1\) \\
 }
 \While{\(r_i \neq 0\)}{
  \(q,\ rem \leftarrow \operatorname{rquot-rem}(r_{i-1}, r_i)\) \\
  \(r_{i+1} \leftarrow rem\) \\
  \(u_{i+1} \leftarrow u_{i-1} - u_i q\) \\
  \(v_{i+1} \leftarrow v_{i-1} - v_{i} q\) \\
  \(i \leftarrow i +1\) \\
 \Return{\(\{u_i, v_i, r_i\}_{i = 0, \dots, h, h+1}\)}
 }
 \caption{Algoritmo extendido de Euclides a la derecha}
\end{algorithm}

\begin{algorithm}[H]\label{alg:left_ext_euc_alg}
 \SetKwInput{KwIn}{Entrada}
 \SetKwInput{KwOut}{Salida}
 \SetKw{Initialization}{Inicialización:}
 \KwIn{\(f,g \in \F[x;\sigma] \text{ con } f \neq 0,\ g \neq 0\)}
 \KwOut{\(\{u_i, v_i, r_i\}_{i = 0, \dots, h, h+1} \text{ tales que } r_i = u_i f + v_i g \text{ para todo  }\) \(0 \le i \le h+1, r_h = (f,g)_r, \text{ y } fu_{h+1} = [f,g]_l.\)}
 \Initialization{\\
 \(r_0 \leftarrow f, r_1 \leftarrow g\) \\
 \(u_0 \leftarrow 1, u_1 \leftarrow 0\) \\
 \(v_0 \leftarrow 0, v_1 \leftarrow 1\) \\
 \(q \leftarrow 0, rem \leftarrow 0\) \\
 \(i \leftarrow 1\) \\
 }
 \While{\(r_i \neq 0\)}{
  \(q,\ rem \leftarrow \operatorname{rquot-rem}(r_{i-1}, r_i)\) \\
  \(r_{i+1} \leftarrow rem\) \\
  \(u_{i+1} \leftarrow u_{i-1} - u_i q\) \\
  \(v_{i+1} \leftarrow v_{i-1} - v_{i} q\) \\
  \(i \leftarrow i +1\) \\
 \Return{\(\{u_i, v_i, r_i\}_{i = 0, \dots, h, h+1}\)}
 }
 \caption{Algoritmo extendido de Euclides a la izquierda}
\end{algorithm}

La prueba del teorema sobre estos algoritmos es una adaptación de la demostración del teorema original para dominios euclideos conmutativos que se encuentra en~\cite{algI}, a excepción del último resultado (el que calcula el mínimo común múltiplo) cuya demostración se encuentra en~\cite{bueso2003algorithmic}.
\begin{theorem}
    Los algoritmos~\ref{alg:right_ext_euc_alg} y~\ref{alg:left_ext_euc_alg} son correctos.
\end{theorem}

\begin{proof}
    Realizaremos únicamente la demostración para el algoritmo~\ref{alg:right_ext_euc_alg}, pues la demostración de la versión a la izquierda es análoga.

    En primer lugar mencionemos que siempre que \(r_i \neq 0\) se tiene que \(\deg(r_{i+1}) < \deg(r_i)\), por tanto existe  \(h \geq 1\) tal que \(r_h \neq 0\) pero \(r_{h+1} = 0\).

Para cada \(i \leq h\) tenemos que \(r_i \neq 0\), y por tanto podemos utilizar la división a la derecha de \(r_{i-1}\) entre \(r_i\) para obtener

\[
r_{i-1} = r_i q_{i+1} + r_{i+1}
.\]

De aquí obtenemos que los divisores a la izquierda comunes de \(r_{i-1}\) y de \(r_i\) coinciden con los divisores a la izquierda comunes de  \(r_i\) y de  \(r_{i+1}\). Luego
\[
r_h = (0, r_h)_l = (r_{h+1}, r_h)_l = (r_h, r_{h-1-})_l = \cdots = (r_1, r_0)_l = (g,f)_l
.\]

A continuación vamos a definir \(u_i, v_i \in R\) con \(i = 0,1, \dots, h, h+1\). En primer lugar, tomamos
\[
u_0 = 1,\ v_0 = 0,\ u_1 = 0,\ v_1 = 1
.\]
Entonces, para \(1 \leq i \leq h\) definimos
\[
u_{i+1} = u_{i-1} - u_iq_{i+1}, \quad
v_{i+1} = v_{i-1} - v_iq_{i+1}
.\]

Aplicaremos un argumento por inducción para comprobar que \(r_i = fu_i + gv_i\). Es inmediato comprobar que para \(i = 0,1\) se cumple, así que supongamos que se cumple para \(i-1,\ i\) y veamos que se cumple para \(i+1\). En efecto
\[
\begin{aligned}
f u_{i+1} + g v_{i+1} =& f(u_{i-1} - u_iq_{i+1}) + g(v_{i-1} - v_iq_{i+1}) = \\
&f u_{i-1} + g v_{i-1} - (f u_i + g v_i)q_{i+1} = r_{i-1} - r_iq_{i+1} = r_{i+1}
\end{aligned}
.\]

Para concluir veamos que \(u_{h+1}f = [f,g]_r\). Observemos en primer lugar que \(0 = r_{h+1} = fu_{h+1} + gv_{h+1}\), luego \(fu_{h+1} = -gv_{h+1}\) es un múltiplo a la derecha común de \(f\) y \(g\). Supongamos \(fu' = -gv'\) un múltiplo común a la derecha de \(f\) y \(g\) cualquiera. Entonces, definimos
\[
\begin{aligned}
    c_0 &= -v' \\
    c_1 &= u' \\
    &\ldots \\
    c_{i+1} &= c_{i-1} - q_{i+1}c_i\quad \text{para}\quad 1 \le i \le h  \\
\end{aligned}
.\]
Así, las siguientes igualdades se cumplen
\[
\begin{aligned}
    r_{i+1}c_{i} - r_{i}c_{i+1} &= 0 \\
    u_{i+1}c_{i} - u_{i}c_{i+1} &= (-1)^{i+1}u' \\
    v_{i+1}c_{i} - v_{i}c_{i+1} &= (-1)^{i+1}v' \\
\end{aligned}
.\]
para \(1 \le i \le h\). Para demostrarlas realizamos una inducción sencilla. Para \(i = 0\), tenemos que
\[
\begin{aligned}
    r_{1}c_{0} - r_{0}c_{1} &= -gv' - fu' =  0 \\
    u_{1}c_{0} - u_{0}c_{1} &= -c_{1} =  -u' \\
    v_{1}c_{0} - v_{0}c_{1} &= c_{0} = -v' \\
\end{aligned}
.\]

Supuestas las tres igualdades ciertas para \(i-1\),
 \[
\begin{aligned}
    r_{i+1}c_{i} - r_{i}c_{i+1} &= (r_{i-1} - r_{i}q_{i+1})c_i - r_i(c_{i-1} - q_{i+1}c_i) = 0 \\
    u_{i+1}c_{i} - u_{i}c_{i+1} &= (u_{i-1} - u_{i}q_{i+1})c_i - u_i(c_{i-1} - q_{i+1}c_i) = u_{i-1}c_i - u_ic_{i-1}= -((-1)^{i}u')  \\
    v_{i+1}c_{i} - v_{i}c_{i+1} &= (v_{i-1} - v_{i}q_{i+1})c_i - v_i(c_{i-1} - q_{i+1}c_i) = v_{i-1}c_i - v_ic_{i-1}= -((-1)^{i}v')  \\
\end{aligned}
.\]
Ahora, por \(r_{h+1} = 0 \neq r_{h}\), sabemos que \(c_{h+1} = 0\), y por tanto \(u_{h+1}\) y \(v_{h+1}\) dividen a la izquierda a \(u'\) y \(v'\) respectivamente. Luego,
\[
f u_{h+1} = -g v_{h+1} = {[f, g]_r}
.\]
\end{proof}

Veamos un lema que nos será útil posteriormente.

\begin{lemma}
\label{lem:reea}
    Sean \(f, g \in \F[x; \sigma] \) y \(  {\{u_{i}, v_{i}, r_{i}\}}_i = 0, \ldots, h\) los coeficientes obtenido al aplicar el Algoritmo Extendido de Euclides a la derecha a \(f\) y \(g\). Notemos \(R_{0} =
    \begin{pmatrix}
    u_0 & u_1 \\
    v_0 & v_1
    \end{pmatrix},\)
    \(Q_i =
    \begin{pmatrix}
        0 & 1 \\
        1 & -q_{i+1}
    \end{pmatrix}
    \)
    y \(R_i =  R_0 Q_1 \cdots Q_i\) para cualquier \(i = 0, \ldots, h\). Por tanto, para todo \(i = 0, \ldots, h\) se cumplen las siguientes afirmaciones:

    \begin{enumerate}
        \item \((fg)R_i = (r_{i} r_{i+1})\).
        \item \(R_i =
            \begin{pmatrix}
            u_i &  u_{i+1} \\
            v_{i} & v_{i+1}
            \end{pmatrix}\).
        \item \(R_i\) tiene inverso a izquierda y derecha.
        \item \((u_i, v_i)_r = 1\).
        \item \(\deg f = \deg r_{i-1} + \deg v_{i}\).

    \end{enumerate}
\end{lemma}
\begin{proof}
Veamos primero las dos primeras afirmaciones, pues probando 2 la afirmación 1 es inmediata por la demostración del algoritmo~\ref{alg:right_ext_euc_alg}. Razonando por inducción, la propiedad es cierta para \(i = 0\), así que supongamos que fijo \(i\) lo es para  \(i-1\). Entonces tenemos que
\[
R_i =
\begin{pmatrix}
    u_{i-1} & u_{i} \\
    v_{i-1} & v_{i}
\end{pmatrix}
\cdot
\begin{pmatrix}
    0 & 1 \\
    1 & -q_{i+1}
\end{pmatrix}
=
\begin{pmatrix}
    u_i & u_{i-1} - u_{i}q_{i+1} \\
    v_i & v_{i-1} - v_{i}q_{i+1} \\
\end{pmatrix}
=
\begin{pmatrix}
    u_{i} & u_{i+1} \\
    v_{i} & v_{i+1}
\end{pmatrix}
,\]

quedando demostrados las afirmaciones 1 y 2. De nuevo si probamos 3 obtenemos 4 inmediatamente. Observemos que \[T_i =
\begin{pmatrix}
q_{i+1} & 1 \\
1 & 0
\end{pmatrix} \]
es una inversa a la izquierda y derecha de \(Q_i\). Por tanto, \(S_i = T_i \cdots T_1 R_0\) es una inversa a izquierda y derecha de \(R_i\) y 3 y 4 quedan probadas.

Veamos 5. Para \(i = 1\),  \(r_0 = f\) y  \(v_1 = 1\) cumpliéndose la igualdad, así que razonemos por inducción. Mencionemos que \(\deg r_i < \deg r_{i-1}\) y que \(\deg v_{i-1} < \deg v_i\) para cualquier \(i > 1\). Entonces, como \(r_{i+1} = r_{i-1} - r_{i}q_{i+1}\) y \(v_{i+1} = v_{i-1} - v_{i}q_{i+1} \) para cualquier \(i\),
 \[
\begin{aligned}
    \deg r_{i-1} = \deg r_{i} + \deg q_{i+1} \\
    \deg v_{i+1} = \deg v_i + \deg q_{i+1}.
\end{aligned}
\]
Ahora, por la hipótesis de inducción, \(\deg f = \deg r_{i-1} + \deg v_{i} = \deg r_{i} + \deg q_{i+1} + \deg v_{i+1} - \deg q_{i+1} = \deg r_i + \deg v_{i+1}\)
\end{proof}

Ahora, con el objetivo de poder definir la evaluación de nuestro polinomios, dado \(j \geq 0\), definimos la \textit{norma j-ésima} para cualquier \(\gamma \in \F\) de forma recursiva como sigue:
\[
N_0(\gamma) = 1
\]
\[
N_{j+1}(\gamma) = \gamma \sigma(N_{j}(\gamma)) = \gamma \sigma(\gamma)\dots\sigma^{j}(\gamma)
.\]
La noción de norma j-ésima admite también una versión para índices negativos dada por
\[
N_{-j-1}(\gamma) = \gamma \sigma^{-1}(N_{-j}(\gamma)) =  \gamma \sigma^{-1}(\gamma) \cdots \sigma^{-j}(\gamma)
.\]

En la siguiente proposición se muestran dos propiedades de esta norma que nos serán de utilidad más adelante.

\begin{proposition}
\label{prop:norm_properties}
Sean \(\gamma, \alpha, \beta \in \F\) tales que \(\beta = \alpha^{-1}\sigma(\alpha)\), entonces
\[
\begin{aligned}
&N_i(\sigma^{k}(\gamma)) = \sigma^{k}(N_i(\gamma)),\\
&N_i(\sigma^k(\beta)) = \sigma^{k}(\alpha)^{-1} \sigma^{k+i}(\alpha).
\end{aligned}\]
\end{proposition}

\begin{proof}
La primera igualdad se prueba de forma inmediata usando la propia definición de la norma, pues
\[
\begin{aligned}
N_i(\sigma^{k}(\gamma)) &= \sigma^{k}(\gamma) \sigma(\sigma^{k}(\gamma))\cdots \sigma^{i-1}(\sigma^{k}(\gamma)) \\
&= \sigma^{k}(\gamma) \sigma^{k}(\sigma(\gamma))\cdots \sigma^{k}(\sigma^{i-1}(\gamma)) = \sigma^{k}(N_i(\gamma)).
\end{aligned}
\]

Para la segunda igualdad, utilizando la igualdad anterior y de nuevo la definición de norma tenemos que
\[
\begin{aligned}
N_i(\sigma^{k}(\beta)) &= \sigma^{k}(N_i(\beta)) = \sigma^{k}\left(\alpha^{-1}\sigma(\alpha)\sigma(\alpha^{-1})\sigma^{2}(\alpha)\sigma^{2}(\alpha^{-1}) \ldots \sigma^{i}(\alpha) \right) \\
   &= \sigma^{k}(\alpha)^{-1} \sigma^{k+i}(\alpha).
\end{aligned}
.\]
\end{proof}

La \textit{evaluación la izquierda} de un polinomio no conmutativo \(f \in R\) en un \(a \in \F\) se define como el resto de la división a la derecha de \(f\) por  \(x - a\), y de forma análoga para la  \textit{evaluación a la derecha}. Estas evaluaciones nos permiten hablar de raíces a izquierda y derecha de estos polinomios. Las propiedades en general de estas son estudiadas en \cite{Lam_1988}, donde se encuentra en esencia la prueba que presentamos del siguiente lema.

\begin{lemma}
\label{lem:eval}
    Sea \(\gamma \in \F\) y  \(f = \sum_{i=0}^n f_ix^i \in R\). Entonces:
    \begin{enumerate}
    \item El resto de la división a la izquerda de \(f\) por  \(x - \gamma\) es  \(\sum_{i=0}^{n} f_i N_i(\gamma)\).
    \item El resto de la división a la derecha de \(f\) por  \(x - \gamma\) es  \(\sum_{i=0}^{n}\sigma^{-i}(f_i)N_{-i}(\gamma)\).
    \item \(N_j(\sigma^k(\gamma)) = \sigma^k(N_j(\gamma))\) para todo \(i,k\).
    \end{enumerate}
\end{lemma}
\begin{proof}
    Para demostrar \textit{i)} observamos primero un caso especial de este resultado:
\begin{equation}
\label{norm_proof}
    x^j - N_j(\gamma) \in R(x-\gamma)\ \forall j \geq 0.
\end{equation}
    Es evidente que el resultado es cierto para \(j = 0\), asi que procedemos por inducción sobre \(j\). Supongamos el resultado cierto para \(j\), entonces
\[
\begin{aligned}
x^{j+1} - N_{j+1}(\gamma) &= x^{j+1} - \sigma(N_j(\gamma))\gamma \\
&= x^{j+1} +  \sigma(N_j(\gamma))(x-\gamma) - \sigma(N_j(\gamma))x \\
&= x^{j+1} +  \sigma(N_j(\gamma))(x-\gamma) - xN_j(\gamma) \\
&= \sigma(N_j(\gamma))(x-\gamma) + x(x^{j} - N_j(\gamma)) \in R(x-\gamma)
\end{aligned}
.\]
Utilizando (\ref{norm_proof}) tenemos entonces que
\[
f - \sum_{i=0}^n f_i N_i(\gamma) = \sum_{i=0}^n f_i(x^i - N_i(\gamma)) \in R(x - \gamma)
\]
y, por la unicidad del resto de la división euclídea tenemos que \(r = \sum_{i=0}^n f_i N_i(\gamma)\).

Para la siguiente afirmación procedemos de forma similar. Veamos en primer lugar que
\begin{equation}
\label{right_norm_proof}
    x^j - N_{-j}(\gamma) \in (x-\gamma)R\ \forall j \geq 0.
\end{equation}

De nuevo, es obvio para \(j = 0\), por tanto procedemos por inducción supuesto cierto para \(j\).
\[
\begin{aligned}
x^{j+1} - N_{-j-1}(\gamma) &= x^{j+1} - \gamma\sigma^{-1}(N_j(\gamma))  \\
&= x^{j+1} +  (x-\gamma)\sigma^{-1}(N_{-j}(\gamma)) - x\sigma^{-1}(N_{-j}(\gamma)) \\
&= x^{j+1} + (x-\gamma)\sigma^{-1}(N_{-j}(\gamma)) - N_{-j}(\gamma)x \\
&= (x-\gamma)\sigma^{-1}(N_{-j}(\gamma)) + (x^{j} - N_{-j}(\gamma))x \in (x-\gamma)R
\end{aligned}
.\]
Así, usando (\ref{right_norm_proof}) vemos que
\[
f - \sum_{i=0}^n \sigma^{-i}(f_i) N_{-i}(\gamma) = \sum_{i=0}^n x^i\sigma^{-i}(f_i) - N_{-i}(\gamma)\sigma^{-i}(f_i)
\]
\[
= \sum_{i=0}^n (x^i - N_{-i}(\gamma))\sigma^{-i}(f_i) \in (x - \gamma)R
.\]
Así que por la unicidad del resto queda probado \textit{ii)}.

Para probar \textit{iii)},
    \[
    N_j(\sigma^k(\gamma)) = \sigma^k(\gamma)\sigma^{k+1}(\gamma)\cdots\sigma^{k+j-1}(\gamma)
\]
\[
    = \sigma^k(\gamma \sigma(\gamma)\dots\sigma^{j-1}(\gamma)) = \sigma^k(N_j(\gamma))
    .\]
\end{proof}

\chapter{Algoritmo para códigos cíclico sesgados}%
\label{chap:conceptos_básicos_sobre_códigos_lineales}

\begin{lemma}
    Sea \(L\) un cuerpo, \(\sigma\) un automorfismo de \(L\) de orden finito \(n\), y \(K = L^\sigma\) el subcuerpo invariante bajo \(\sigma\). Sea  \(\{a_0, \dots, a_{n-1}\) una \(K\)-base de \(L\). Entonces, para todo \(t \leq n\), y cada subconjunto \(k_0 < k_1 < \dots < k_{t-1} \subset \{0, 1, \dots, n-1\}\)
    \[
    \begin{vmatrix}
        \alpha_{k_0} & \alpha_{k_1} & \dots & \alpha_{k_{t -1}} \\
        \sigma^{t-1}(\alpha_{k_0}) & \sigma^{t-1}(\alpha_{k_1}) & \dots & \sigma^{t-1}(\alpha_{k_{t-1}}) \\
        \vdots & \vdots & \ddots & \vdots \\
        \sigma^{t-1}(\alpha_{k_0}) & \sigma^{t-1}(\alpha_{k_1}) & \dots & \sigma^{t-1}(\alpha_{k_{t-1}})
    \end{vmatrix}
    \neq 0
    .\]
\end{lemma}

\begin{proofs}
    Realizaremos la prueba por inducción sobre \(t\). El caso \(t = 1\) se cumple trivialmente. Por tanto, supongamos que el lema se cummple para un cierto  \(t \geq 1\). Tenemos que comprobar que, para toda matriz \((t+1) \times (t+1)\)
    \[
    \Delta =
    \begin{pmatrix}
        \alpha_{k_0} & \alpha_{k_1} & \dots & \alpha_{k_{t -1}} \\
        \sigma^{t-1}(\alpha_{k_0}) & \sigma^{t-1}(\alpha_{k_1}) & \dots & \sigma^{t-1}(\alpha_{k_{t-1}}) \\
        \vdots & \vdots & \ddots & \vdots \\
        \sigma^{t}(\alpha_{k_0}) & \sigma^{t}(\alpha_{k_1}) & \dots & \sigma^{t}(\alpha_{k_{t}})

    \end{pmatrix}
    .\]
el determinante \(|\Delta|\) es distinto de cero. Supongamos por el contrario que \(|\Delta| = 0\). Por la hipótesis de inducción tenemos que las primeras \(t\) columnas de \(\Delta\) son linealmente independientes, luego existen \(a_0, \dots, a_{t-1} \in L\) tales que

\[
(a_{k_t}, \sigma^{t-1}(a_{k_t}, \dots, \sigma^t(a_{k_t}))) = \sum_{j=0}^{t-1} a_j(\alpha_{k_j}, \sigma^{t-1}(\alpha_{k_j}), \dots, \sigma^t(\alpha_{k_j}))
.\]

Es decir, \(a_0, \dots, a_{t-1}\) satisfacen el sistema lineal

\begin{equation}
\label{linear_system}
\begin{cases}
    \alpha_{k_t} = a_0\alpha_{k_0} + \cdots + a_{t-1}\alpha_{k_{t-1}} \\
    \sigma^{t-1}(\alpha_{k_t}) = a_0\sigma^{t-1}(\alpha_{k_0}) + \cdots + a_{t-1}\sigma^{t-1}(\alpha_{k_{t-1}}) \\
    \vdots \\
    \sigma^t(\alpha_{k_t}) = a_0\sigma^t(\alpha_{k_0}) + \cdots + a_{t-1}\sigma^t(\alpha_{k_{t-1}})
\end{cases}
.
\end{equation}

Para cada \(j = 0, \dots, t-1\), restamos en (\ref{linear_system}) la ecuación \(j+1\) transformada por \(\sigma^{-1}\) a la ecuación \(j\). Esto produce el siguiente sistema lineal homogéneo

\begin{equation}
\label{hom_liner_system}
\begin{cases}
0 = (a_0 - \sigma^{-1}(a_0))\alpha_{k_0} + \cdots + (a_{t-1} - \sigma^{-1}(a_{t-1}))\alpha_{k_{t-1}} \\
0 = (a_0 - \sigma^{-1}(a_0))\sigma^{t-1}(\alpha_{k_0}) + \cdots + (a_{t-1} - \sigma^{-1}(a_{t-1}))\sigma^{t-1}(\alpha_{k_{t-1}}) \\
\vdots \\
0 = (a_0 - \sigma^{-1}(a_0))\sigma^{t-1}(\alpha_{k_0}) + \cdots + (a_{t-1} - \sigma^{-1}(a_{t-1}))\sigma^{t-1}(\alpha_{k_{t-1}}) \\
\end{cases}
.
\end{equation}

La matriz de coeficientes de (\ref{hom_liner_system}) es no singular, por la hipótesis de inducción, así que tenemos que para todo \(j = 0, \dots, t-1\), \(a_j - \sigma^{-1}(a_j) = 0\), y por tanto \(a_0, \dots, a_{t-1} \in K\). Como consecuencia, la ecuación (\ref{linear_system}) establece una dependencia lineal sobre  \(K\) de la  \(K\)-base  \(\{\alpha_0, \dots, \alpha{n-1}\), creando una contradicción. Por tanto, \(|\Delta| \neq 0\) y el resultado queda demostrado.
\end{proofs}

\begin{lemma}
    Sea \(\alpha \in \F\) tal que \(\{\alpha, \sigma(\alpha), \dots, \sigma^{n-1}(\alpha)\) sea una base de \(\F\) como  \(\F^\sigma\)-espacio vectorial. Fijemos  \(\beta = \alpha^{-1}\sigma(\alpha)\). Para todo subconjunto \(T = \{t_1 < t_2 < \cdots < t_m\} \subset \{0, 1, \dots, n-1\}\), los polinomios
    \[
    g^l = [x - \sigma^{t_1}(\beta), x - \sigma^{t_2}(\beta), \dots, x - \sigma^{t_m}(\beta)]_l
    \]
y
    \[
    g^r = [x - \sigma^{t_1}(\beta^{-1}), x - \sigma^{t_2}(\beta^{-1}), \dots, x - \sigma^{t_m}(\beta^{-1})]_r
    \]
tienen grado \(m\). Por tanto, si \(x - \sigma^s(\beta) |_r g^l\) o \(x - \sigma^s(\beta^{-1}) |_l g^r\), entonces \(s \in T\).
\end{lemma}

\begin{proofs}
    Supongamos que \(\deg g^l < m\), así que \(g^l = \sum_{i=0}^{m-1} g_i x^i\). Como \(g\) es un múltiplo a la izquierda de \(x - \sigma^{t_j}(\beta)\) para todo \(1 \leq j \leq m\), por el lema (TODO) tenemos que
    \[
    \sum_{i=0}^{m-1}g_i N_i(\sigma^{t_j}(\beta)) = 0 \text{ para todo } 1 \leq j \leq m
    .\]

Esto es un sistema lineal homogéneo cuya matriz de coeficientes es la transpuesta de
\end{proofs}

\chapter{Implementación en SageMath del Algoritmo de Sugiyama para códigos cíclicos sesgados}%
\label{chap:implementación_en_sagemath_del_algoritmo_de_sugiyama_para_códigos_cíclicos_sesgados}

Se ha desarrollado una implementación en SageMath del algoritmo de Sugiyama para códigos cíclicos sesgados, junto con la estructura necesaria para representar estos códigos.

Para ello hemos seguido las directrices de SageMath presentes en
\begin{center}
\url{https://doc.sagemath.org/html/en/thematic_tutorials/structures_in_coding_theory.html}.
\end{center}
 Para realizar dicha implementación hemos utilizado varias clases abstractas de Sage cuya finalidad es servir de base para implementar nuevas familias de códigos, junto con los algoritmos asociados de codificación y decodificación. Además Sage contiene una implementación de anillos de polinomios sesgados, la cual también se usa a lo largo de todo el desarrollo del algoritmo y de las clases para represetnar la estructura.

La primera clase que se ha implementado es \texttt{SkewCyclicCode} para representar cualquier código cíclico sesgado, la cual hereda de la clase abstracta \texttt{AbstractLinearCode}. Para poder codificar mensajes en palabras código de este tipo de códigos se han diseñado las clases \texttt{SkewCyclicCodeVectorEncoder} y \texttt{SkewCyclicCodePolynomialEncoder}, ambas heredando de la clase de Sage \texttt{Encoder}. A continuación, heredando de la clase \texttt{SkewCyclicCode} se implementó la clase \texttt{SkewRSCode} que almacena la estructura necesaria para trabajar con códigos Reed-Solomon sesgados. Por último, la clase \texttt{SkewRSCodeSugiyamaDecoder}, que hereda de la clase \texttt{Decoder}, contiene la implementación del algoritmo de decodificación desarrollado en este trabajo.

A continuación mostramos la documentación junto con ejemplos de uso de cada una de las estructuras y funciones implementadas.

\newpage
\section{Clase de códigos cíclicos sesgados}%
\label{sec:clase_de_códigos_cíclicos_sesgados}

\begin{description}[font=\ttfamily, style=nextline]
    \item[class SkewCyclicCode(self, generator\_pol=None)]

    Clase para representar un código cíclico sesgado. Hereda de \texttt{AbstractLinearCode}.

    \textsc{Argumentos}
    \begin{description}[font=\ttfamily, style=nextline]
        \item[generator\_pol] Polinomio generador para construir el código.
    \end{description}

    \textsc{Ejemplos}
    \begin{lstlisting}
sage: F.<t> = GF(3^10)
sage: sigma = F.frobenius_endomorphism()
sage: R.<x> = F['x', sigma]
sage: g = x**5 - 1
sage: C = SkewCyclicCode(generator_pol=g)
sage: C
[10, 5] Skew Cyclic Code over Finite Field in
t of size 3^10
    \end{lstlisting}

    \textsc{Métodos}
    \begin{description}[font=\ttfamily, style=nextline]
    \item[generator\_pol(self)]
    Devuelve el polinomio generador de \texttt{self}.

    \textsc{Ejemplos}
    \begin{lstlisting}
sage: F.<t> = GF(3^10)
sage: sigma = F.frobenius_endomorphism()
sage: R.<x> = F['x', sigma]
sage: g = x**5 - 1
sage: C = SkewCyclicCode(generator_pol=g)
sage: C.generator_polynomial()
x^5 + 2
    \end{lstlisting}

    \end{description}

    \begin{description}[font=\ttfamily, style=nextline]
    \item[skew\_polynomial\_ring(self)]
    Devuelve el anillo de polinomios de sesgados de \texttt{self}.

    \textsc{Ejemplos}
    \begin{lstlisting}
sage: F.<t> = GF(3^10)
sage: sigma = F.frobenius_endomorphism()
sage: R.<x> = F['x', sigma]
sage: g = x**5 - 1
sage: C = SkewCyclicCode(generator_pol=g)
sage: C.generator_polynomial()
Skew Polynomial Ring in x over Finite Field in t of
size 3^10 twisted by t |--> t^3
    \end{lstlisting}

    \end{description}

\end{description}

\section{Clases de codificadores para códigos cíclicos sesgados}

\begin{description}[font=\ttfamily, style=nextline]
    \item[class SkewCyclicCodeVectorEncoder(self, code)]
    Clase para codificar un vector como palabras código. Hereda de \texttt{Encoder}.

    \textsc{Argumentos}

    \begin{description}[font=\ttfamily, style=nextline]
        \item[code]
        Código sobre el que codificaremos los vectores.

    \end{description}

    \textsc{Ejemplos}

    \begin{lstlisting}
sage: F.<t> = GF(3^10)
sage: sigma = F.frobenius_endomorphism()
sage: R.<x> = F['x', sigma]
sage: g = x**5 - 1
sage: C = SkewCyclicCode(generator_pol=g)
sage: E = SkewCyclicCodeVectorEncoder(C)
sage: E
Vector-style encoder for [10, 5] Skew Cyclic
Code over Finite Field in t of size 3^10
    \end{lstlisting}

    \textsc{Métodos}

    \begin{description}[font=\ttfamily, style=nextline]
        \item[generator\_matrix(self)]
        Construye una matriz generatriz asociada al código asociado a este codificador

        \textsc{Ejemplos}
    \begin{lstlisting}
sage: F.<t> = GF(3^10)
sage: sigma = F.frobenius_endomorphism()
sage: R.<x> = F['x', sigma]
sage: g = x**5 - 1
sage: C = SkewCyclicCode(generator_pol=g)
sage: E = SkewCyclicCodeVectorEncoder(C)
sage: E.generator_matrix()
[2 0 0 0 0 1 0 0 0 0]
[0 2 0 0 0 0 1 0 0 0]
[0 0 2 0 0 0 0 1 0 0]
[0 0 0 2 0 0 0 0 1 0]
[0 0 0 0 2 0 0 0 0 1]
    \end{lstlisting}
    \end{description}

\end{description}

\begin{description}[font=\ttfamily, style=nextline]
    \item[class SkewCyclicCodePolynomialEncoder(self, code)]
    Clase para codificar polinomios en palabras código. Hereda de \texttt{Encoder}.

    \textsc{Argumentos}
    \begin{description}[font=\ttfamily, style=nextline]
        \item[code]
        Código sobre el que codificaremos los vectores.
    \end{description}

    \textsc{Ejemplos}
    \begin{lstlisting}
sage: F.<t> = GF(3^10)
sage: sigma = F.frobenius_endomorphism()
sage: R.<x> = F['x', sigma]
sage: g = x**5 - 1
sage: C = SkewCyclicCode(generator_pol=g)
sage: E = SkewCyclicCodePolynomialEncoder(C)
sage: E
Polynomial-style encoder for [10, 5] Skew Cyclic
Code over Finite Field in t of size 3^10
    \end{lstlisting}

    \textsc{Métodos}
    \begin{description}[font=\ttfamily, style=nextline]
        \item[encode(self, p)]
        Transforma el polinomio \texttt{p} en un elemento del código asociado a \texttt{self}.

        \textsc{Argumentos}
        \begin{description}[font=\ttfamily, style=nextline]
            \item[p]
            Polinomio que será codificado.
        \end{description}

        \textsc{Ejemplos}
        \begin{lstlisting}
sage: F.<t> = GF(3^10)
sage: sigma = F.frobenius_endomorphism()
sage: R.<x> = F['x', sigma]
sage: g = x**5 - 1
sage: C = SkewCyclicCode(generator_pol=g)
sage: E = SkewCyclicCodePolynomialEncoder(C)
sage: m = x + t + 1
sage: E.encode(m)
(2*t + 2, 2, 0, 0, 0, t + 1, 1, 0, 0, 0)
        \end{lstlisting}

        \item[unencode\_nocheck(self, c)]
        Devuelve el mensaje en forma de polinomio asociado a \texttt{c} sin comprobar si \texttt{c} pertenece al código

        \textsc{Argumentos}
        \begin{description}[font=\ttfamily, style=nextline]
            \item[c]
            Un vector de la misma longitud del código asociado a \texttt{self}
        \end{description}

        \textsc{Ejemplos}
        \begin{lstlisting}
sage: F.<t> = GF(3^10)
sage: sigma = F.frobenius_endomorphism()
sage: R.<x> = F['x', sigma]
sage: g = x**5 - 1
sage: C = SkewCyclicCode(generator_pol=g)
sage: E = SkewCyclicCodePolynomialEncoder(C)
sage: m = x + t + 1
sage: w = E.encode(m)
sage: E.unencode_nocheck(w)
x + t + 1
        \end{lstlisting}

        \item[message\_space(self)]
        Devuelve el espacio de mensajes asociado a \texttt{self}.

        \textsc{Ejemplos}
        \begin{lstlisting}
sage: F.<t> = GF(3^10)
sage: sigma = F.frobenius_endomorphism()
sage: R.<x> = F['x', sigma]
sage: g = x**5 - 1
sage: C = SkewCyclicCode(generator_pol=g)
sage: E = SkewCyclicCodePolynomialEncoder(C)
sage: E.message_space()
Skew Polynomial Ring in x over Finite Field in t
of size 3^10 twisted by t |--> t^3
        \end{lstlisting}
    \end{description}

\end{description}

\section{Clase de Códigos Reed-Solomon Sesgados}

\begin{description}[font=\ttfamily, style=nextline]
    \item[class SkewRSCode(self, hamming\_dist=None, \\skew\_polynomial\_ring=None, alpha=None)]

    Clase para representar un código RS sesgado. Hereda de \texttt{SkewCyclicCode}.

    \textsc{Argumentos}
    \begin{description}[font=\ttfamily, style=nextline]
        \item[hamming\_dist] La distancia hamming de \texttt{self}.
        \item[skew\_polynomial\_ring] El anillo de polinomios sesgados base.
        \item[alpha] Un generador normal de la extensión de cuerpos sobre el cuerpo fijo por el automorfismo asociado a \texttt{skew\_polynomial\_ring}.
    \end{description}

    \textsc{Ejemplos}
    \begin{lstlisting}
sage: F.<t> = GF(2^12, modulus=x**12 + x**7 +
              x**6 + x**5 + x**3 + x +1)
sage: sigma = (F.frobenius_endomorphism())**10
sage: R.<x> = F['x', sigma]
sage: alpha = t
sage: RS_C = SkewRSCode(hamming_dist=5,
              skew_polynomial_ring=R, alpha=alpha)
sage: RS_C
[6, 2] Skew Reed Solomon Code over Finite
Field in t of size 2^12
    \end{lstlisting}

    \textsc{Métodos}
    \begin{description}[font=\ttfamily, style=nextline]
        \item[hamming\_dist(self)]
            Devuelve la distancia Hamming asociada a este código.

        \textsc{Ejemplos}
        \begin{lstlisting}
sage: F.<t> = GF(2^12, modulus=x**12 + x**7 +
              x**6 + x**5 + x**3 + x +1)
sage: sigma = (F.frobenius_endomorphism())**10
sage: R.<x> = F['x', sigma]
sage: alpha = t
sage: RS_C = SkewRSCode(hamming_dist=5,
              skew_polynomial_ring=R, alpha=alpha)
sage: RS_C.hamming_dist()
2
        \end{lstlisting}
    \end{description}
\end{description}

\section{Clase del decodificador para códigos RS sesgados}
\begin{description}[font=\ttfamily, style=nextline]
    \item[class SkewRSCodeSugiyamaDecoder(self, code)]
    Clase para decodificar un vector recibido en una palabra del código asociado a \texttt{self}. Hereda de \texttt{Decoder}.

    \textsc{Argumentos}
    \begin{description}[font=\ttfamily, style=nextline]
        \item[code]
        Código al que pertenecerá la palabra que devuelve el decodificador.
    \end{description}

    \textsc{Ejemplos}
    \begin{lstlisting}
sage: F.<t> = GF(2^12, modulus=x**12 + x**7 +
              x**6 + x**5 + x**3 + x +1)
sage: sigma = (F.frobenius_endomorphism())**10
sage: R.<x> = F['x', sigma]
sage: alpha = t
sage: RS_C = SkewRSCode(hamming_dist=5,
              skew_polynomial_ring=R, alpha=alpha)
sage: D = SkewRSCodeSugiyamaDecoder(RS_C)
sage: D
Decoder through the Sugiyama like algorithm of the
[6, 2] Skew Reed Solomon Code over Finite Field in t
of size 2^12
    \end{lstlisting}

    \textsc{Métodos}
    \begin{description}[font=\ttfamily, style=nextline]
        \item[decode\_to\_code(self, word)]
        Decodifica el vector \texttt{word} en un elemento del código asociado a \texttt{self}.

        \textsc{Argumentos}
        \begin{description}[font=\ttfamily, style=nextline]
            \item[word]
            Vector de elemtos del mismo cuerpo sobre el que está construido el anillo de polinomios
            sesgados del código asociado a \texttt{self}.
        \end{description}

        \textsc{Ejemplos}
        \begin{lstlisting}
sage: F.<t> = GF(2^12, modulus=x**12 + x**7 +
              x**6 + x**5 + x**3 + x +1)
sage: sigma = (F.frobenius_endomorphism())**10
sage: R.<x> = F['x', sigma]
sage: alpha = t
sage: RS_C = SkewRSCode(hamming_dist=5,
              skew_polynomial_ring=R, alpha=alpha)
sage: D = SkewRSCodeSugiyamaDecoder(RS_C)
sage: P_E = SkewCyclicCodePolynomialEncoder(RS_C)
sage: m = x + t
sage: codeword = P_E.encode(m)
sage: noisy_codeword = copy(codeword)
sage: noisy_codeword[3] = t**671
sage: decoded_word = D.decode_to_code(noisy_codeword)
sage: codeword == decoded_word
True
        \end{lstlisting}

        \item[decoding\_radius(self)]
        Devuelve el número de errores que \texttt{self} es capaz de corregir.

        \textsc{Ejemplos}
        \begin{lstlisting}
sage: F.<t> = GF(2^12, modulus=x**12 + x**7 +
              x**6 + x**5 + x**3 + x +1)
sage: sigma = (F.frobenius_endomorphism())**10
sage: R.<x> = F['x', sigma]
sage: alpha = t
sage: RS_C = SkewRSCode(hamming_dist=5,
              skew_polynomial_ring=R, alpha=alpha)
sage: D = SkewRSCodeSugiyamaDecoder(RS_C)
sage: D.decoding_radius()
2
        \end{lstlisting}
    \end{description}
\end{description}

\section{Funciones auxiliares}

\begin{description}[font=\ttfamily, style=nextline]
    \item[left\_extended\_euclidean\_algorithm(skew\_polynomial\_ring, f, g)]

    Implementación del algoritmo extendido de Euclides a la izquierda.

    \textsc{Argumentos}
    \begin{description}[font=\ttfamily]
        \item[skew\_polynomial\_ring] anillo de polinomios sesgados
        \item[f] polinomio del anillo pasado como argumento
        \item[g] polinomio del anillo pasado como argumento de grado menor que \texttt{f}
    \end{description}

    \textsc{Salida}
    \begin{description}[font=\ttfamily]
        \item[u, v r] vectores de polinomios sesgados tales que \texttt{u[i]*f + v[i]*g = r[i]} para cualquier posición \texttt{i}
    \end{description}

    \textsc{Ejemplos}
    \begin{lstlisting}
sage: F.<t> = GF(3^10)
sage: sigma = F.frobenius_endomorphism()
sage: R.<x> = F['x', sigma]
sage: a = x + t +1;
sage: b = t^2 * x**2 + (t+1)*x +1
sage: f = a*b
sage: g = b
sage: left_extended_euclidean_algorithm(R, f, g)
([1, 0, 1],
 [0, 1, 2*x + 2*t + 2],
 [t^6*x^3 + (2*t^3 + t^2 + 1)*x^2 +
    (t^2 + 2*t + 2)*x + t + 1, t^2*x^2 +
    (t + 1)*x + 1, 0])
    \end{lstlisting}

    \item[right\_extended\_euclidean\_algorithm(skew\_polynomial\_ring, f, g)]

    Implementación del algoritmo extendido de Euclides a la derecha.

    \textsc{Argumentos}
    \begin{description}[font=\ttfamily]
        \item[skew\_polynomial\_ring] anillo de polinomios sesgados
        \item[f] polinomio del anillo pasado como argumento
        \item[g] polinomio del anillo pasado como argumento de grado menor que \texttt{f}
    \end{description}

    \textsc{Salida}
    \begin{description}[font=\ttfamily]
        \item[u, v r] vectores de polinomios sesgados tales que \texttt{f*u[i] + g*v[i] = r[i]} para cualquier posición \texttt{i}
    \end{description}

    \textsc{Ejemplos}
    \begin{lstlisting}
sage: F.<t> = GF(3^10)
sage: sigma = F.frobenius_endomorphism()
sage: R.<x> = F['x', sigma]
sage: a = x + t +1;
sage: b = t^2 * x**2 + (t+1)*x +1
sage: f = a*b
sage: g = a
sage: right_extended_euclidean_algorithm(R, f, g)
([1, 0, 1],
 [0, 1, 2*t^2*x^2 + (2*t + 2)*x + 2],
 [t^6*x^3 + (2*t^3 + t^2 + 1)*x^2 +
    (t^2 + 2*t + 2)*x + t + 1, x + t + 1, 0])
    \end{lstlisting}

    \item[left\_lcm(pols)]
    Calcula el mínimo común múltiplo a la izquierda para una lista de polinomios sesgados

    \textsc{Argumentos}
    \begin{description}[font=\ttfamily, style=nextline]
        \item[pols] Lista de polinomios sesgados
    \end{description}

    \textsc{Ejemplos}
    \begin{lstlisting}
sage: F.<t> = GF(3^10)
sage: sigma = F.frobenius_endomorphism()
sage: R.<x> = F['x', sigma]
sage: a = x + t +1;
sage: b = t^2 * x**2 + (t+1)*x +1
sage: left_lcm((a,b))
x^3 +
(2*t^9 + 2*t^8 + 2*t^7 + t^6 + t^5 + 2*t^4 + t^2 + t)*x^2 +
(2*t^9 + t^8 + 2*t^7 + 2*t^6 + 2*t^5 + 2*t^3 + t)*x +
t^8 + t^6 + 2*t^4 + t^3 + 2*t + 2
    \end{lstlisting}

    \item[norm(j, sigma, gamma)]
    Calcula la norma \texttt{j}-ésima de \texttt{gamma}.

    \textsc{Argumentos}
    \begin{description}[font=\ttfamily]
        \item[j] Número entero
        \item[sigma] Automorfismo del cuerpo donde esté definido \texttt{gamma}
        \item[gamma] Un elemento de un cuerpo cualquiera
    \end{description}

    \textsc{Ejemplos}
    \begin{lstlisting}
sage: F.<t> = GF(3^10)
sage: sigma = F.frobenius_endomorphism()
sage: R.<x> = F['x', sigma]
sage: norm(2, sigma, t**2 + t)
t^8 + t^7 + t^5 + t^4
    \end{lstlisting}

\end{description}

% Añadir tantos capítulos como sea necesario

% \cleardoublepage\part{Segunda parte}
% \input{capitulos/capitulo02}

% --------------------------------------------------------------------
% APPENDIX: Opcional
% --------------------------------------------------------------------

% \appendix % Reinicia la numeración de los capítulos y usa letras para numerarlos
% \pdfbookmark[-1]{Apéndices}{appendix} % Alternativamente podemos agrupar los apéndices con un nuevo \part{Apéndices}

% \input{apendices/apendice01}
% Añadir tantos apéndices como sea necesario

% --------------------------------------------------------------------
% GLOSARIO: Opcional
% --------------------------------------------------------------------

% % !TeX root = ../libro.tex
% !TeX encoding = utf8

\chapter*{Glosario}
\addcontentsline{toc}{chapter}{Glosario} % Añade el glosario a la tabla de contenidos

La inclusión de un glosario es opcional.

Archivo: \texttt{glosario.tex}

\begin{description} 
  \item[$\mathbb{R}$] Conjunto de números reales.

  \item[$\mathbb{C}$] Conjunto de números complejos.

  \item[$\mathbb{Z}$] Conjunto de números enteros.
\end{description}
\endinput


% -------------------------------------------------------------------
% BACKMATTER
% -------------------------------------------------------------------

\backmatter % Desactiva la numeración de los capítulos
\pdfbookmark[-1]{Referencias e Índices}{BM-Referencias}

% BIBLIOGRAFÍA
%-------------------------------------------------------------------

\setbibpreamble{Las referencias se listan por orden alfabético. Aquellas referencias con más de un autor están ordenadas de acuerdo con el primer autor.\par\bigskip}
\bibliographystyle{alpha}
\begin{small} % Normalmente la bibliografía se imprime en un tamaño de letra más pequeño.
\bibliography{library.bib}
\end{small}

% ÍNDICE TERMINOLÓGICO  (Opcional)
%-------------------------------------------------------------------

\cleardoublepage
\begin{footnotesize} % Normalmente el índice se imprime en un tamaño de letra más pequeño.
\printindex
\end{footnotesize}

\end{document}
